\chapter{形式化验证}\label{ch:formal-verification}

\section{形式化验证简介}\label{sec:introduction-to-formal-verification}

嵌入式实时操作系统的本质是复杂性、多样性和可靠性。操作系统的硬件无关部分包括大约2000行代码，各个部分之间具有较复杂的作用和联系方式。作为一个操作系统又要应付不同应用的各种需求，并且需要保证功能的正确性。同时，我们还需要提高开发效率，控制成本和开发周期，以及保证软件的顺利交付。

传统的软件设计方法基于自然语言的思考、设计和描述，往往片面和模糊，极易引起误解。它也无法进行严格的检查，只能通过人的心智进行分析。基于
等半形式化的方法采用一些相对清晰的图形化描述，一些工具也能自动生成代码框架并检查分析。以上两种方法在测试系统时，均是设计一系列用例对其进行测试，最多有结构化测试的参与。但是，它们都无法保证系统中没有错误，均不适用于性命攸关的系统的分析和开发，因此我们需要更为严格的开发设计流程。完全的形式化方法则基于严格定义的数学概念和语言，可以开发自动化工具进行检查和分析。它把数学的严格性带入软件开发的各个阶段，通过严格的数学证明保证系统中没有漏洞。

软件安全性的标准分为7个等级\hyperref[sec:chapter-bibliography-8]{{[}1{]}}，分别如下\footnote{也列入在内}：

\captionof{table}{EAL等级及其描述}\label{tab:eal-levels-and-their-descriptions}

{\def\LTcaptype{none} % do not increment counter
\begin{longtable}[]{@{}
  >{\raggedright\arraybackslash}p{(\linewidth - 2\tabcolsep) * \real{0.1175}}
  >{\raggedright\arraybackslash}p{(\linewidth - 2\tabcolsep) * \real{0.8825}}@{}}
\toprule\noalign{}
\endhead
\bottomrule\noalign{}
\endlastfoot
EAL级别 & 描述 \\
EAL1 &
EAL1是经过功能性测试的系统。它提供基本的安全保障目标描述。这些描述中包含了系统的功能规范和潜在错误分析。在EAL1中，每一个安全保障目标都要列出并且评估测试。

相比于未评估系统，EAL1仅仅是一个小规模的安全提升。 \\
EAL2 &
EAL2是使用基本软件工程方法论进行结构化开发的系统。它提供全面的安全保障目标描述。这些描述中包含了系统的基本功能规范和基本安全架构。在EAL2中，每一个安全保障目标都要列出并且进行独立的评估测试；此外还需要有一定攻击能力的安全测试人员对这些安全保障目标进行攻击，以确定没有问题。

EAL2相比于EAL1有更严格的测试规范和系统弱点分析。 \\
EAL3 &
EAL3是使用强化的软件工程方法论进行系统性测试和检查的系统。它提供全面的安全保障目标描述。这些描述中包含了系统的基本功能规范和安全架构的详细设计。在EAL3中，每一个安全保障目标都要列出并且进行独立的评估测试；此外还需要有一定攻击能力的安全测试人员针对安全架构的详细设计进行攻击，以确定没有问题。

EAL3相比于EAL2有更严格的系统架构规范和相应的测试。 \\
EAL4 &
EAL4是使用强化的软件工程方法论进行系统性设计、测试和评估的系统。它提供全面的安全保障目标描述。这些描述中包含了系统的详细功能规范和安全架构的详细设计，还要包含一部分关键模块的划分和具体实现。在EAL4中，每一个安全保障目标都要列出并且进行独立评估测试；此外还需要有高于平均水平攻击能力的安全测试人员针对安全架构的详细设计进行攻击，以确定没有问题。

EAL4相比于EAL3有更严格的设计描述、实现审查和更强大的测试。 \\
EAL5 &
EAL5是使用半形式化方法进行设计和测试的系统。它提供全面的安全保障目标描述。这些描述中包含了系统的详细功能规范和安全架构的详细设计，还要包含半形式化方法描述的全部关键模块的划分和一部分模块的具体实现。在EAL5中，每一个安全保障目标都要列出并且进行独立评估测试；此外还需要有相当程度攻击能力的安全测试人员针对安全架构的详细设计进行攻击，以确定没有问题。

EAL5相比于EAL4有半形式化的设计描述、更清晰的架构和更强大的测试。 \\
EAL6 &
EAL6是使用半形式化方法进行设计、测试和验证的系统。它提供全面的安全保障目标描述。这些描述中包含了系统的详细功能规范和安全架构的详细设计，还要包含全部关键模块的划分和全部模块的具体实现。EAL6要求对于安全保障目标的策略有完整的形式化描述，对其功能则要有半形式化描述。具体的功能设计必须是模块化的、分层的和简明的。在EAL6中，每一个安全保障目标都要列出并且进行独立评估测试；此外还需要有较强攻击能力的安全测试人员针对安全架构的详细设计进行攻击，以确定没有问题。

EAL6相比于EAL5有更全面的设计分析和弱点测试、以及结构化的实现描述。 \\
EAL7 &
EAL7是形式化方法进行设计、测试和验证的系统。它提供全面的安全保障目标描述。这些描述中包含了系统的详细功能规范和安全架构的详细设计，还要包含全部关键模块的划分和全部模块的结构化的具体实现。EAL7要求对于安全保障目标的策略有完整的形式化描述，对其功能则要有半形式化描述。具体的功能设计必须是模块化的、分层的和简明的。在EAL7中，每一个安全保障目标都要列出并且针对实现进行完整的独立评估测试；此外还需要有极强攻击能力的安全测试人员针对安全架构的详细设计进行攻击，以确定没有问题。

EAL7相比于EAL6有更进一步的形式化设计和更全面的测试。 \\
EAL7+ &
EAL7+是严格形式化方法进行设计、测试和验证的系统。它提供全面的安全保障目标描述。这些描述中包含了系统的详细功能规范和安全架构的详细设计，还要包含全部关键模块的划分和全部模块的结构化的具体实现。EAL7+要求对于安全保障目标的功能和策略有完整的形式化描述。具体的功能设计必须是模块化的、分层的和简明的，并且要证明具体的实现完全符合形式化描述本身。在EAL7+中，每一个安全保障目标都要列出并且针对实现进行完整的独立评估测试，测试用例必须由形式化工具自动生成；此外还需要有极强攻击能力的安全测试人员针对安全架构的详细设计进行攻击，以确定没有问题。

EAL7+相比于EAL7有更进一步的形式化实现和形式化测试。

RMP系统是采用EAL7+标准设计和测试的。 \\
\end{longtable}
}

绝大多数操作系统如FreeRTOS、Windows、Linux的认证级别都在EAL4。少数几款其他系统的认证在EAL5，INTEGRITY-178B的认证级别在EAL6+。由于RMP的结构相对简单，而且其安全目标中只包括功能正确性目标而不包括信息安全目标\footnote{多数//等设计规范不适用}，因此较为容易取得更高的认证水平。

\section{系统的形式化规范}\label{sec:formal-specification-of-the-system}

在中，由于架构相关部分非常短小\footnote{往往在50行代码以内}，容易保证其正确性，因此不对该部分进行形式化验证。另一个原因是架构相关部分往往随芯片而发生变化，而要支持的平台多达数十，从人力成本上讲不方便对其一一建立形式化模型进行验证。因此，主要的验证工作集中在文件的1500行代码上。注释不计算在内，约有1000行之多。基于同样的多架构支持原因，对这些代码的正确性，我们仅验证至语言语义的层面，并不进行二进制指令级别的验证。这意味着，最终可执行文件的正确性仍然需要依赖于编译器的正确性。由于在实践中编译器通常具备较高的认证等级，因此信任编译器的做法是合理的。当然，如果不希望信任编译器，也可以采用经形式化验证的\hyperref[sec:chapter-bibliography-8]{{[}2{]}}编译器。

形式化验证并不能够完全杜绝软硬件系统错误操作的发生，尤其是硬件错误导致的问题\footnote{比如CPU或外设的硅片错误，这些错误的描述常常可以在各个厂商的芯片勘误手册内找到}。对于硬件错误和硬件抽象层（Hardware
Abstraction
Layer，HAL）实现错误导致的系统失灵，的形式化验证是无能为力的。

\section{形式化证明}\label{sec:formal-proof}

将完成的后续工作。

\section{其他文档}\label{sec:other-documents}

将完成的后续工作。

\section{本章参考文献}\label{sec:chapter-bibliography-7}

\protect\phantomsection\label{sec:chapter-bibliography-8}{[}1{]}
Common Criteria. "Common Criteria for Information Technology Security
Evaluation", Part 3: Security assurance components, 2012.

{[}2{]} X. Leroy, "The CompCert verified compiler," Documentation and
user manual. INRIA Paris-Rocquencourt, 2012.
